\documentclass[11pt]{article}
\usepackage[letterpaper,margin=1in]{geometry}
\usepackage{amsmath,amssymb,amsthm,booktabs,tabularx,microtype,xurl}
\usepackage[hidelinks]{hyperref}
\newtheorem{theorem}{Theorem}
\newtheorem{lemma}[theorem]{Lemma}
\newtheorem{proposition}[theorem]{Proposition}
\newcommand{\C}{\mathbb C}
\newcommand{\Q}{\mathbb Q}
\newcommand{\Mon}{\operatorname{Mon}}
\newcommand{\disc}{\operatorname{disc}}

\title{Full wreath-product monodromy for the square\\
of an explicit Keller map}
\author{Alec Kriebel\\[2pt]\large with heavy assistance from ChatGPT 5.6 Sol}
\date{Research draft prepared 21 July 2026, 17:00:46 UTC\\
Not peer reviewed or publicly released as a numbered exploration}

\begin{document}
\maketitle

\begin{abstract}
Let $F:\C^3\to\C^3$ be the newly announced noninjective polynomial map with
constant Jacobian determinant $-2$ and geometric monodromy $S_3$. We determine
the geometric monodromy of its canonical self-composition:
\[
  \Mon(F\circ F)=S_3\wr S_3
\]
in degree nine, of order $1296$. On the target line $(1,2,s)$, an exact
degree-nine eliminant has a one-edge Newton polygon at infinity, giving a
9-cycle; its discriminant has a squarefree factor of multiplicity one, giving
a single within-block transposition; and the outer cubic gives a transposition
of the three blocks. An elementary group lemma forces the full wreath product.
Postcomposition by $\operatorname{diag}(1/4,1,1)$ gives a noninjective
unit-Jacobian map realizing the pair $(9,S_3\wr S_3)$. We also prove that every
iterate $F^{\circ m}$ has a full $3^m$-cycle in its geometric inertia at
infinity. Exact SymPy, dependency-free finite-field, PARI/GP, and GAP checks
accompany the argument.
\end{abstract}

\noindent\textbf{Verification disclaimer.}
The human author is a complete amateur exploring the limits of AI-assisted
mathematics and cannot independently verify these claims. This is a research
artifact for expert checking, not an established result. No guarantee of
worldwide priority is made.

\section{The map and the result}

Put $u=1+xy$ and define $F=(F_1,F_2,F_3)$ by
\begin{align}
F_1&=u^3z+y^2u(4+3xy),\nonumber\\
F_2&=y+3xu^2z+3xy^2(4+3xy),\label{eq:F}\\
F_3&=2x-3x^2y-x^3z.\nonumber
\end{align}
Exact expansion gives $\det JF=-2$. Moreover,
\begin{equation}
p_0=(0,0,-\tfrac14),\quad
p_1=(1,-\tfrac32,\tfrac{13}2),\quad
p_2=(-1,\tfrac32,\tfrac{13}2)
\label{eq:points}
\end{equation}
all map to $(-1/4,0,0)$.

For a dominant polynomial map $H$, write $\Mon(H)$ for the Galois group of
the normal closure of the induced function-field extension, in its natural
action on a generic fiber.

\begin{theorem}\label{thm:main}
The self-composition $G=F\circ F$ has generic degree nine and
\[
\Mon(G)=S_3\wr S_3\leq S_9,
\qquad |\Mon(G)|=6^3\cdot6=1296.
\]
The action has three blocks of size three, indexed by the intermediate points
in a generic fiber of the outer copy of $F$.
\end{theorem}

\begin{proposition}\label{prop:iterates}
For every $m\geq1$, the geometric monodromy of $F^{\circ m}$ contains a cycle
of length $3^m$.
\end{proposition}

\begin{proposition}\label{prop:normalized}
Let $\widehat G=(G_1/4,G_2,G_3)$. Then $\det J\widehat G=1$,
$\Mon(\widehat G)=S_3\wr S_3$, and $\widehat G$ is noninjective. The three
points in~\eqref{eq:points} all map under $\widehat G$ to $(0,0,-1/2)$.
\end{proposition}

The last proposition follows immediately from the theorem and
$\det J(F\circ F)=4$.

\section{The cubic inverse resolvent}

For a target $(a,b,c)$, introduce
\begin{equation}
C_{a,b,c}(t)=2at^3-bt^2+2t-c.
\label{eq:cubic}
\end{equation}
For a simple root $t$, the generic preimage is reconstructed by
\begin{equation}
y=-\frac{bt^2+3c-6t}{2t^2},\qquad
x=\frac{t}{1-ty},\qquad
z=\frac{2x-3x^2y-c}{x^3}.
\label{eq:reconstruct}
\end{equation}
Substitution reduces every coordinate of $F(x,y,z)-(a,b,c)$ to zero modulo
$C_{a,b,c}(t)$. The parameter is $t=x/(1+xy)$, reciprocal to the usual
parameter $y+1/x$.

The three roots of~\eqref{eq:cubic} parametrize the generic fiber. For
$F\circ F$, each of three outer roots supports a second cubic of inner roots.
Consequently the generic degree is nine and
\begin{equation}
\Mon(F\circ F)\leq S_3\wr S_3.
\label{eq:upper}
\end{equation}

\section{A degree-nine slice}

Restrict the target to $(a,b,c)=(1,2,s)$. Let $t$ satisfy
\[
2t^3-2t^2+2t-s=0
\]
and reconstruct $(x(t),y(t),z(t))$ using~\eqref{eq:reconstruct}. An inner
resolvent root $r$ satisfies
\[
2x(t)r^3-y(t)r^2+2r-z(t)=0.
\]
Eliminating $t$ and removing the vertical content $256s^7$ gives the primitive
polynomial $P(r,s)$ below.

\begingroup\small
\begin{align*}
P={}&128r^9s-256r^8
+(2592s^2-3072s+1408)r^7\\
&+(-7452s^3+1248s^2-1584s-576)r^6\\
&+(-17496s^4+82944s^3-87936s^2+42240s-5760)r^5\\
&+(52488s^4-142884s^3+135144s^2-60144s+8928)r^4\\
&+(104976s^5+97200s^4-97056s^3-41472s^2+93824s-26880)r^3\\
&+(236196s^6-380538s^5-147744s^4+772080s^3-754944s^2\\
&\hspace{34mm}+333792s-62208)r^2\\
&+(236196s^6-734832s^5+1916784s^4-2658432s^3+2128320s^2\\
&\hspace{34mm}-919296s+186624)r\\
&+531441s^8-2204496s^7+5436882s^6-11805534s^5+17996446s^4\\
&\hspace{34mm}-17402304s^3+10364976s^2-3544992s+557280.
\end{align*}
\endgroup

Its roots generically parametrize the nine points in the slice fiber.

\section{Three inertia elements}

\subsection{A 9-cycle at infinity}

Write $v=1/s$ and consider $v^8P(r,v^{-1})$. For the coefficients of
$r^0,\ldots,r^9$, their $v$-adic valuations are
\[
0,2,2,3,4,4,5,6,8,7.
\]
Every intermediate point lies strictly above the line from $(0,0)$ to $(9,7)$.
The lower Newton polygon is therefore one edge of slope $7/9$.

Set $v=w^9$ and $r=w^{-7}R$. The initial equation is
\[
128R^9+531441=0.
\]
Its roots are simple and lift uniquely to Puiseux branches. The local deck
transformation $w\mapsto\zeta_9w$ multiplies the leading term by
$\zeta_9^{-7}$. Since $\gcd(7,9)=1$, it permutes all nine branches in one
orbit. Thus slice monodromy contains a 9-cycle $\alpha$.

\subsection{A single transposition}

Exact discriminant factorization gives
\begin{equation}
\disc_r(P)=2^{38}q(s)^8A(s)B(s)^2,
\qquad q(s)=27s^2-28s+12,
\label{eq:disc}
\end{equation}
where
\begingroup\small
\begin{align*}
A(s)={}&1129718145924s^{12}-6474958632657s^{11}
+21955119111630s^{10}\\
&-60661410535386s^9+123515279853390s^8-191093865073182s^7\\
&+235490291194248s^6-188631015819696s^5+54561577345568s^4\\
&+41882989175328s^3-44714997684096s^2+15094447332352s\\
&-1434819245056.
\end{align*}
\endgroup
The exact certificates verify that $A$ is squarefree and coprime to $q$, $B$,
and the leading coefficient $128s$. At each root of $A$, the discriminant has
valuation one while the degree remains nine. Local geometric inertia is thus a
single transposition $\beta$. A permutation moving only two sheets cannot act
nontrivially on the three blocks, so $\beta$ is supported in one block.

The full degree-22 factor $B$ is recorded in Appendix~\ref{app:B} and checked
independently in SymPy and PARI/GP.

\subsection{A transposition of the blocks}

The outer cubic has discriminant
\[
\disc_t(2t^3-2t^2+2t-s)=-4(27s^2-28s+12)=-4q(s).
\]
The quadratic $q$ is squarefree and coprime to $A$ and $B$. A small loop about
either root of $q$ has outer monodromy a transposition. Total inertia contains
an element $\gamma$ whose image on the three blocks is that transposition.

\section{The group argument}

\begin{lemma}\label{lem:group}
Let $H\leq S_3\wr S_3$ in its natural action on three blocks of size three.
Suppose $H$ contains a 9-cycle $\alpha$, a single transposition $\beta$, and
an element $\gamma$ inducing a transposition on the blocks. Then
$H=S_3\wr S_3$.
\end{lemma}

\begin{proof}
The 9-cycle induces a 3-cycle on the blocks, and $\alpha^3$ acts as a 3-cycle
inside each block. The transposition $\beta$ is supported in one block.
Conjugating it by $\alpha^3$ gives a distinct transposition in that block, so
these elements generate the full $S_3$ supported there. Conjugation by powers
of $\alpha$ gives the full $S_3$ independently in all three blocks. Hence the
base subgroup $S_3^3$ lies in $H$.

On blocks, $\alpha$ supplies a 3-cycle and $\gamma$ a transposition, so the
image of $H$ is the top $S_3$. A subgroup containing the full base and
surjecting onto the top is the full semidirect product.
\end{proof}

The three inertia elements constructed above satisfy the lemma, so the slice
monodromy is $S_3\wr S_3$. Slice monodromy is a subgroup of generic monodromy;
combining this with~\eqref{eq:upper} proves Theorem~\ref{thm:main}.

\section{Full-cycle inertia for all iterates}

We prove Proposition~\ref{prop:iterates}. Put $v=1/s$, start from
$X_0=(1,2,s)$, and choose a compatible inverse tower
$F(X_k)=X_{k-1}$. Write
\[
x_k\sim c_xs^{-a_k},\qquad y_k\sim c_ys^{a_k},\qquad
z_k\sim c_zs^{c_k},
\]
where the constants are nonzero. Initially $a_0=0$ and $c_0=1$.

The next inverse parameter satisfies
\[
2x_kt^3-y_kt^2+2t-z_k=0.
\]
If $c_k>5a_k$, the lower Newton edge balances the first and last terms, so
\[
t_{k+1}\sim c_ts^{(a_k+c_k)/3}.
\]
The reconstruction formulas then yield
\begin{equation}
a_{k+1}=\frac{c_k-2a_k}{3},\qquad
c_{k+1}=2(c_k-a_k).
\label{eq:recurrence}
\end{equation}
The hypothesis persists: for $\rho_k=a_k/c_k$, the interval
$0\leq\rho_k\leq1/6$ is invariant under
\[
\rho_{k+1}=\frac{1-2\rho_k}{6(1-\rho_k)},
\]
whose image on that interval is $[2/15,1/6]$ after the first step. Thus no
omitted term lies on the lower edge and no leading coefficient cancels.

Set $A_k=3^ka_k$ and $C_k=3^kc_k$. Equation~\eqref{eq:recurrence} becomes
\[
A_{k+1}=C_k-2A_k,\qquad C_{k+1}=6(C_k-A_k),
\quad A_0=0,\ C_0=1.
\]
The exponent of $t_m$ is
\[
\frac{E_m}{3^m},\qquad E_m=A_{m-1}+C_{m-1}.
\]
For $m\geq2$, $C_{m-1}$ is divisible by three and $A_{m-1}\equiv1\pmod3$.
Hence $3\nmid E_m$. The first exponents are
\[
\frac13,\quad\frac79,\quad\frac{34}{27},\quad\frac{178}{81},\ldots.
\]
At level $m$, the local inverse field has ramification index at least $3^m$;
the global degree is at most $3^m$, so equality holds and the local extension
is totally ramified. Finite extensions of $\C((v))$ are tame Puiseux
extensions, so inertia acts as a single $3^m$-cycle. This proves the
proposition.

As an independent check, nested PARI resultants for $m=3$ give a degree-27
eliminant with one lower Newton edge from $(0,0)$ to $(27,34)$.

\section{Independent arithmetic cross-check}

At $s=-3$, the primitive specialization is
\begingroup\small
\begin{align*}
f(r)={}&-384r^9-256r^8+33952r^7+216612r^6-4580568r^5\\
&+9515052r^4-15697056r^3+223986114r^2+599887620r\\
&+17172300267.
\end{align*}
\endgroup
Dependency-free modular arithmetic verifies squarefree factorizations with
irreducible-degree patterns
\[
\begin{array}{c|c}
p&\text{degrees}\\ \hline
13&(9)\\
61&(2,1,1,1,1,1,1,1)\\
19&(2,2,2,2,1).
\end{array}
\]
These give Frobenius cycle types $(9)$, $(2,1^7)$, and $(2^4,1)$. The same
group lemma forces the arithmetic group to be the full wreath product. GAP
independently enumerates all subgroups of the wreath product and finds no
proper survivor. PARI/GP 2.17.4 identifies five unrelated rational fibers as
transitive group $9T31$ of order $1296$.

This is corroboration only; the geometric proof does not infer geometric
monodromy from an arithmetic specialization.

\section{Verification and priority scope}

The artifact contains four independent checks:
\begin{center}
\begin{tabularx}{\textwidth}{@{}lX@{}}
\toprule
File & Scope\\ \midrule
\texttt{verify\_symbolic.py} & Jacobian, collision, inverse identities,
resultant, Newton data, discriminant, squarefreeness, coprimality\\
\texttt{verify\_modular.py} & dependency-free polynomial arithmetic and Rabin
irreducibility certificates\\
\texttt{verify\_pari.gp} & independent discriminant, coprimality, finite-field,
and arithmetic Galois checks\\
\texttt{verify\_group.g} & exhaustive subgroup check inside $S_3\wr S_3$\\
\bottomrule
\end{tabularx}
\end{center}
The optional \texttt{verify\_level3\_newton.py} checks the degree-27 edge with
SymPy-generated equations and PARI resultants.

The closest public monodromy study states only
$\Mon(F_0\circ F_0)\leq S_3\wr S_3$ and calls its exact group a natural next
computation~\cite{weighted}. A separate four-variable construction already
claims a different imprimitive degree-eight group of order $192$~\cite{quartic}.
We therefore do not claim the first imprimitive or first non-symmetric Keller
monodromy. The claimed contribution is the exact canonical self-composition
group, the realization $(9,S_3\wr S_3)$, and the all-iterate full-cycle
statement. A timestamped 23-repository audit accompanies the artifact.

The stronger arboreal assertion that every iterate has the full iterated
wreath product remains open: the present argument does not yet produce a new
single transposition at every level.

\appendix
\section{The degree-22 square factor}\label{app:B}
\begingroup\scriptsize
\begin{align*}
B(s)={}&1423119505038213888s^{22}-155058052818404824443s^{21}
+1413754646027656083066s^{20}\\
&-8145658955220494785812s^{19}+33677018812011807334224s^{18}
-108385682498649366622416s^{17}\\
&+280859820926917245978240s^{16}-598631883156564470992728s^{15}
+1062288637590844067815584s^{14}\\
&-1579910926841521168581900s^{13}+1974730823883289142626824s^{12}
-2073348574060015592229024s^{11}\\
&+1822485349587628391880960s^{10}-1333523184375693801555072s^9
+806144943610022071172864s^8\\
&-398941760796329492797440s^7+159756878618286560778240s^6
-50952869532774134959104s^5\\
&+12636860641575849510912s^4-2344566242453066735616s^3
+304677812001210531840s^2\\
&-24551525690267664384s+926711257485017088.
\end{align*}
\endgroup

\begin{thebibliography}{9}
\bibitem{mo}
\emph{Galois structure of the new counterexample to the Jacobian conjecture},
MathOverflow 513387 (2026),
\url{https://mathoverflow.net/questions/513387/}.

\bibitem{weighted}
MikhailSzh and Claude,
\emph{The Antiderivative Resolvent of the Weighted-Lift Family of Keller Maps},
audited Git commit \texttt{9193c66385ce390f61a4e25d3f5255435bfa056a} (2026),
\url{https://github.com/MikhailSzh/weighted-lift-galois}.

\bibitem{quartic}
Juan M. G. H.,
\emph{A quartic-resultant counterexample to the Jacobian Conjecture},
audited Git commit \texttt{640322133fa4bbe172e0ac95f3c485f2c86d8cea} (2026),
\url{https://github.com/juanmgh3/quartic-resultant-keller-map}.

\bibitem{composition}
J.~K\"onig, D.~Neftin, and S.~Rosenberg,
\emph{Polynomial compositions with large monodromy groups and applications to
arithmetic dynamics}, arXiv:2401.17872 (2024).
\end{thebibliography}

\end{document}
